% !TEX root = ../main.tex
\section{Conclusions}

This Master's Thesis has presented the motivation, theoretical basis and development strategy for \texttt{tpeanuts}, a modular, publicly available PyTorch framework for neutrino-oscillation phenomenology that grows out of, and substantially extends, the PEANUTS toolkit. The project is based on a conservative principle for the part it inherits directly from PEANUTS: its physical conventions must be preserved, while the numerical implementation is modernised to support tensor operations, batching and accelerator execution; everything built on top of that foundation is new.

That foundation is the organisation of the PEANUTS numerical core into a PyTorch-oriented design covering PMNS construction, matter potentials, vacuum propagation, matter mixing and solar-neutrino probabilities, together with a validation methodology based on low-level comparisons, intermediate physical quantities and final probabilities. This methodology is essential because speed is useful only if the updated implementation remains physically faithful, and it is the precondition for everything that follows, not the main result in itself.

The central contribution of this work is the modular software architecture described in \cref{sec:pipelines,sec:validation-strategy}, which decouples the core oscillation engine from the specific neutrino source, matter-density model and physics scenario being evaluated. This design has made it possible to extend the scope of the original PEANUTS code, limited to solar-neutrino propagation, to atmospheric-neutrino detection through the Earth; to support multiple interchangeable solar and Earth density models and alternative perturbative treatments of the matter potential; and to incorporate beyond-the-Standard-Model extensions, namely non-standard interactions and sterile neutrinos, without modifying the underlying propagation engine. The validation and benchmark results of \cref{sec:validation-strategy,sec:experimental-reproduction-results,sec:bsm-results} confirm that these extensions integrate correctly with the rest of the framework and reproduce established physical signatures, demonstrating that the modular design is not only a software-engineering convenience but a working foundation for future physics studies.

A second, equally central contribution is the extension of this engine into a real, data-driven statistical-inference framework (\cref{sec:detector-layer,sec:inference-layer}): a detector-agnostic forward-folding layer and a gradient-based likelihood/fitting layer, built directly on PyTorch autograd, connect the same validated oscillation engine to real, published data from five experiments. The Daya Bay case study (\cref{sec:dayabay-case-study}) demonstrates that this chain works end to end, a near/far ratio fit recovering both \(\sin^22\theta_{13}\) and \(\Delta m^2_{31}\) jointly within \(1\)--\(2\%\) of Daya Bay's published values, and its diagnosed limitations (the full-spectral-shape fit's non-convergence without a full systematic covariance, the Laplace approximation's breakdown for SNO~II's \(\Delta m^2_{21}\)) identify concrete next steps rather than open questions.

Taken together, these two contributions, the modular oscillation engine and the detector-inference layer built on top of it, are released publicly (\cref{sec:supplementary-material}) as a documented, installable Python package rather than kept as a closed thesis codebase. Combined with its batched, GPU-capable and differentiable design, this is intended to make \texttt{tpeanuts} directly reusable by other researchers working on solar, atmospheric or BSM neutrino phenomenology, not only a record of the results reproduced in this thesis.

Future work will focus on five lines: making the remaining detector-level analyses (full-shape reactor fit, Borexino's background model, SNO~I, IceCube's free systematics) as robust as the Daya Bay near/far and SNO~II fits already are, mainly through richer correlated systematic treatments; deepening the BSM extensions, in particular connecting the sterile extension to the real-data inference machinery already validated for the Standard-Model detector layer; extending the current Majorana-neutrino parametrisation, presently limited to the two extra CP phases that leave every oscillation probability exactly unchanged, with genuine magnetic-field physics, a transition magnetic moment and the resulting spin-flavour precession Hamiltonian, the channel through which Majorana and Dirac neutrinos actually become experimentally distinguishable; generalising the propagation medium towards an arbitrary-density-profile formulation; and combining the individually validated detector-inference chains into a genuine joint likelihood across experiments.
